\subsubsection{Synthetic Eddy Method (SEM) Boundary Condition}%
\label{SEC:BC:Inflow:SEM}

The Synthetic Eddy Method (SEM) provides an unsteady turbulent inflow boundary
condition that is crucial for accurate LES without requiring large run-up
sections of computational grid for fully developing a turbulent boundary layer
prior to the computational geometry of interest.
%
The user provides an inflow profile with the desired time-averaged flow
properties and turbulence statistics.
%
At each time step, the SEM will then generate an unsteady inflow profile with
turbulent fluctuations that, when time averaged, will reproduce the requested
distribution.
%
\Sref{SEC:SEM} provides extensive details on the synthetic eddy method:
%
\sref{SEC:SEM:theory} provides the theory behind the method as well as an
overview of how the method works;
%
\sref{SEC:SEM:usage} provides details on how to use the SEM in GFR including the
file format for the SEM profile, the relevant input variables with some
recommendations, and other aspects unique to the SEM implementation within GFR;
%
and \sref{SEC:SEM:FAQ} provides a collection of frequently asked questions.

%###############################################################################

